\documentclass[11pt]{article}

\usepackage[T1]{fontenc}
\usepackage{amsmath,amssymb,amsthm}
\usepackage{hyperref}

\title{Consistent Communality Rules for the Guttman Estimator}
\author{}
\date{\today}

\newcommand{\T}{^{\mathsf T}}
\newcommand{\E}{\mathbb{E}}
\newcommand{\Var}{\operatorname{Var}}
\newcommand{\Cov}{\operatorname{Cov}}
\newcommand{\tr}{\operatorname{tr}}
\newcommand{\diag}{\operatorname{diag}}
\newcommand{\rank}{\operatorname{rank}}
\newcommand{\vech}{\operatorname{vech}}
\newcommand{\argmin}{\operatorname*{arg\,min}}

\theoremstyle{plain}
\newtheorem{proposition}{Proposition}
\newtheorem{lemma}{Lemma}
\theoremstyle{remark}
\newtheorem*{remark}{Remark}

\begin{document}
\maketitle

\section*{Purpose}

Dhaene and Rosseel \cite{dhaene2024} define their multiple-group estimator by
replacing the diagonal of the sample covariance matrix with estimated
communalities. They use Spearman's rule as described by Harman
\cite{harman1976}, and note that other communality or uniqueness estimators
exist, including the uniqueness estimator of Ihara and Kano \cite{iharakano1986}
and Cudeck \cite{cudeck1991}. This note records the population requirement a
communality rule must satisfy before it can be used as part of the Guttman
estimator rather than merely as a start-value heuristic.

The required target is the diagonal of the common covariance,
\[
  c_i(\Sigma_0)
  =
  e_i\T \Lambda\Phi\Lambda\T e_i,
\]
not an arbitrary lower bound, reliability coefficient, or squared multiple
correlation. In a simple-structure CFA block, where indicator \(i\) loads only
on factor \(f(i)\),
\[
  c_i=\lambda_i^2\Phi_{f(i)f(i)}.
\]

\section{Population setup}

Consider one factor block \(I_f\) with at least three indicators and diagonal
residual covariance within the block:
\begin{equation}\label{eq:block-model}
  y_i=\lambda_i\eta_f+\epsilon_i,\qquad
  \Var(\eta_f)=\phi_f,\qquad
  \Cov(\epsilon_i,\epsilon_j)=0\quad(i\ne j).
\end{equation}
For \(i\ne j\) in the same block,
\[
  \sigma_{ij}=\lambda_i\lambda_j\phi_f.
\]
The common diagonal entry for item \(i\) is
\[
  c_i=\lambda_i^2\phi_f,
  \qquad
  \theta_i=\sigma_{ii}-c_i.
\]
The Guttman matrix \(C\) is obtained by replacing \(\sigma_{ii}\) by \(c_i\), or
equivalently by subtracting \(\theta_i\) from the diagonal.

\begin{remark}
For a multi-factor simple CFA, the same within-factor formulas hold factor by
factor. Cross-factor covariances involve \(\Phi_{fg}\) and should not be used in
the communality triads below. The target remains
\(\diag(\Lambda\Phi\Lambda\T)\).
\end{remark}

\section{Spearman triad identities}

\begin{proposition}[Single triad]\label{prop:triad}
Let \(i,j,k\in I_f\) be distinct indicators with
\(\lambda_j\lambda_k\phi_f\ne0\). Under \eqref{eq:block-model},
\begin{equation}\label{eq:single-triad}
  h_{i;jk}(\Sigma)
  =
  \frac{\sigma_{ij}\sigma_{ik}}{\sigma_{jk}}
  =
  c_i .
\end{equation}
\end{proposition}

\begin{proof}
Substitute the within-factor covariances:
\[
  \frac{\sigma_{ij}\sigma_{ik}}{\sigma_{jk}}
  =
  \frac{(\lambda_i\lambda_j\phi_f)(\lambda_i\lambda_k\phi_f)}
       {\lambda_j\lambda_k\phi_f}
  =
  \lambda_i^2\phi_f
  =
  c_i .
\]
\end{proof}

Thus every admissible pair \((j,k)\) gives a population-consistent estimator of
the same diagonal common variance. Pooling pairs changes finite-sample behavior,
not the population target.

\section{Two consistent Spearman poolings}

Let
\[
  \mathcal P_i=\{(j,k): j<k,\ j,k\in I_{f(i)},\ j\ne i,\ k\ne i\}.
\]
Assume all required denominators are nonzero.

\paragraph{Average of ratios.}
The lavaan helper used by \texttt{lav\_cfa\_theta\_spearman()} and the current
\texttt{magmaan} helper averages the triad ratios on the correlation scale. On
the covariance scale the unclamped population functional is
\begin{equation}\label{eq:avg-ratio}
  h_i^{\mathrm{AR}}(\Sigma)
  =
  \frac{1}{|\mathcal P_i|}
  \sum_{(j,k)\in\mathcal P_i}
  \frac{\sigma_{ij}\sigma_{ik}}{\sigma_{jk}} .
\end{equation}
By Proposition~\ref{prop:triad}, \(h_i^{\mathrm{AR}}(\Sigma_0)=c_i\).

\paragraph{Ratio of sums.}
Dhaene and Rosseel's printed formula is the Harman/Spearman ratio-of-sums
pooling:
\begin{equation}\label{eq:ratio-sums}
  h_i^{\mathrm{RS}}(\Sigma)
  =
  \frac{\sum_{(j,k)\in\mathcal P_i}\sigma_{ij}\sigma_{ik}}
       {\sum_{(j,k)\in\mathcal P_i}\sigma_{jk}} .
\end{equation}
Under \eqref{eq:block-model},
\[
  \sum_{(j,k)}\sigma_{ij}\sigma_{ik}
  =
  \lambda_i^2\phi_f
  \sum_{(j,k)}\sigma_{jk},
\]
so \(h_i^{\mathrm{RS}}(\Sigma_0)=c_i\) whenever the denominator does not vanish.

\begin{remark}
The two poolings are equivalent on the exact model but differ off model and in
finite samples. The ratio of sums is often numerically calmer when individual
\(\sigma_{jk}\)'s are small. The average of ratios avoids cancellation in the
pooled denominator but can be sensitive to a single weak pair. Both are genuine
population-consistent communality rules under the simple block model.
\end{remark}

\section{Linear method-of-moments triad family}

The ratio identity can be written without ratios. Work on the correlation
scale within one factor block. Under \eqref{eq:block-model},
\[
  \rho_{ij}=\alpha_i\alpha_j\qquad(i\ne j),
\]
where
\[
  \alpha_i=\lambda_i\sqrt{\phi_f/\sigma_{ii}} .
\]
The correlation-scale communality proportion is
\[
  \gamma_i=\alpha_i^2,
  \qquad
  c_i=\sigma_{ii}\gamma_i .
\]
For each pair \((j,k)\in\mathcal P_i\), the Spearman identity is equivalent to
the linear moment
\begin{equation}\label{eq:linear-triad-moment}
  \gamma_i\rho_{jk}-\rho_{ij}\rho_{ik}=0 .
\end{equation}
Thus Spearman's rule is not an ad hoc population identity. It is one way of
pooling exact rank-one triad moments. The estimator-design question is which
finite-sample pooling of \eqref{eq:linear-triad-moment} to use.

For a sample correlation matrix \(R\), define for item \(i\)
\[
  a_{i,jk}=R_{jk},
  \qquad
  b_{i,jk}=R_{ij}R_{ik},
  \qquad
  (j,k)\in\mathcal P_i .
\]
Stack these into vectors \(a_i\) and \(b_i\). For any symmetric positive
semidefinite weight \(W_i\) satisfying \(a_i\T W_i a_i>0\), define
\begin{equation}\label{eq:item-linear-mom}
  \hat\gamma_i(W_i)
  =
  (a_i\T W_i a_i)^{-1}a_i\T W_i b_i,
  \qquad
  \hat c_i(W_i)=S_{ii}\hat\gamma_i(W_i).
\end{equation}
This is the closed-form weighted least-squares estimator of \(\gamma_i\) in
the linear moments
\[
  a_i\gamma_i-b_i=0 .
\]

\begin{proposition}[Linear triad moments are exact]\label{prop:linear-mom}
Let \(\Sigma_0\) satisfy the one-factor block model
\eqref{eq:block-model}. Let \(R_0\) be its correlation matrix, and suppose
\(W_i(R_0)\) is finite with
\[
  a_i(R_0)\T W_i(R_0)a_i(R_0)>0 .
\]
Then the population version of \eqref{eq:item-linear-mom} satisfies
\[
  \gamma_i\{W_i(R_0)\}=\alpha_i^2,
  \qquad
  c_i\{W_i(R_0)\}=\lambda_i^2\phi_f .
\]
\end{proposition}

\begin{proof}
At the population model,
\[
  b_{i,jk}(R_0)
  =
  \rho_{ij}\rho_{ik}
  =
  \alpha_i^2\rho_{jk}
  =
  \gamma_i a_{i,jk}(R_0)
\]
for every \((j,k)\in\mathcal P_i\). Hence \(b_i(R_0)=a_i(R_0)\gamma_i\), and
\[
  (a_i\T W_i a_i)^{-1}a_i\T W_i b_i
  =
  (a_i\T W_i a_i)^{-1}a_i\T W_i a_i\,\gamma_i
  =
  \gamma_i .
\]
Multiplying by \(\sigma_{ii}\) returns \(c_i=\lambda_i^2\phi_f\).
\end{proof}

Several familiar and useful rules are special cases.

\paragraph{Identity linear moments.}
With \(W_i=I\),
\begin{equation}\label{eq:identity-linear-mom}
  \hat\gamma_i^{\mathrm{ILM}}
  =
  \frac{\sum_{(j,k)\in\mathcal P_i}R_{jk}R_{ij}R_{ik}}
       {\sum_{(j,k)\in\mathcal P_i}R_{jk}^2}.
\end{equation}
This is the simplest closed-form method-of-moments version. It avoids division
by each individual \(R_{jk}\), so it is less sensitive than the average of
ratios to a single weak denominator.

\paragraph{Average-of-ratios Spearman.}
If all denominators are nonzero, the average-of-ratios rule
\eqref{eq:avg-ratio} is obtained from \eqref{eq:item-linear-mom} by the
data-dependent weight
\[
  W_i=\diag\{R_{jk}^{-2}:(j,k)\in\mathcal P_i\}.
\]
Then \(a_i\T W_i a_i=|\mathcal P_i|\) and
\[
  \hat\gamma_i(W_i)
  =
  |\mathcal P_i|^{-1}
  \sum_{(j,k)\in\mathcal P_i}
  \frac{R_{ij}R_{ik}}{R_{jk}} .
\]

\paragraph{Ratio-of-sums Spearman.}
When the relevant \(R_{jk}\)'s have compatible sign, the ratio-of-sums rule
\eqref{eq:ratio-sums} is the same linear estimating equation with row weights
proportional to \(R_{jk}^{-1}\):
\[
  \hat\gamma_i^{\mathrm{RS}}
  =
  \frac{\sum_{(j,k)\in\mathcal P_i}R_{ij}R_{ik}}
       {\sum_{(j,k)\in\mathcal P_i}R_{jk}} .
\]
This form is consistent whenever the pooled denominator has a nonzero
population limit, but it is not always a positive least-squares metric if the
within-block covariances have mixed signs.

\section{One-step GMM weights for the triad moments}

The itemwise construction can be stacked. Let \(\gamma\) collect the
correlation-scale communality proportions for all indicators. For every
admissible row \((i;j,k)\), write
\[
  g_{i;jk}(\gamma;R)
  =
  \gamma_iR_{jk}-R_{ij}R_{ik}.
\]
Stacking rows gives
\[
  g(\gamma;R)=A(R)\gamma-b(R),
\]
where each row of \(A(R)\) has a single nonzero entry \(R_{jk}\) in column
\(i\), and the corresponding row of \(b(R)\) is \(R_{ij}R_{ik}\). For any
positive semidefinite weight \(W\),
\begin{equation}\label{eq:stacked-gmm}
  \hat\gamma_W
  =
  \{A(R)\T W A(R)\}^{+}
  A(R)\T W b(R).
\end{equation}
This is a closed-form GMM estimator for the overidentified linear triad
moments. The identity choice \(W=I\) gives a stable baseline. A one-step
efficient version uses a preliminary consistent \(\tilde\gamma\), forms
\[
  \widehat\Omega_g
  =
  \widehat{\operatorname{AVar}}\{
    \sqrt n\,g(\tilde\gamma;R)
  \},
\]
and sets \(W=\widehat\Omega_g^{+}\) in \eqref{eq:stacked-gmm}. The covariance
\(\widehat\Omega_g\) is obtained by delta method from the asymptotic covariance
of the sample correlations; under normality this has a closed form, and under
nonnormality it can be estimated by a sandwich covariance of the correlation
vector.

\begin{remark}
The one-step estimator is efficient only within the declared triad-moment GMM
family. It is not a claim that the resulting full Guttman estimator is globally
efficient for CFA. The weight also affects finite-sample behavior and the
off-model diagonal-completion target.
\end{remark}

This gives a recommendable hierarchy. The average-of-ratios and ratio-of-sums
rules are historically recognizable certified rules. The identity linear-moment
rule is the simplest non-ratio method-of-moments default. The one-step GMM rule
is the natural efficiency-oriented member of the same closed-form family,
conditional on a chosen covariance model for the sample correlations.

\section{Uniqueness-first rules}

The paper also points to uniqueness estimators \cite{iharakano1986,cudeck1991}.
For the Guttman estimator, the relevant population object is not the uniqueness
itself but
\[
  h_i=\sigma_{ii}-u_i .
\]
A uniqueness-first rule is acceptable only if it is exact on the population
factor model.

\begin{proposition}[Exact uniqueness implies exact communality]
\label{prop:uniqueness}
Suppose a population functional \(u(\Sigma)\in\mathbb R^p\) is defined on a
class of identified factor models and satisfies
\[
  u(\Sigma_0)=\diag(\Theta_0)
\]
whenever
\[
  \Sigma_0=\Lambda_0\Phi_0\Lambda_0\T+\Theta_0
\]
with diagonal \(\Theta_0\). Then
\[
  h_i^u(\Sigma)=\sigma_{ii}-u_i(\Sigma)
\]
satisfies \(h_i^u(\Sigma_0)=e_i\T\Lambda_0\Phi_0\Lambda_0\T e_i\).
\end{proposition}

\begin{proof}
At \(\Sigma_0\),
\[
  \sigma_{0,ii}
  =
  e_i\T\Lambda_0\Phi_0\Lambda_0\T e_i+\Theta_{0,ii}.
\]
Subtracting \(u_i(\Sigma_0)=\Theta_{0,ii}\) gives the claim.
\end{proof}

One exact population reading is the rank-\(k\) uniqueness problem: find a
diagonal \(\Theta\) such that \(\Sigma-\Theta\) is positive semidefinite, has
rank \(k\), and satisfies the required loading-pattern restrictions. If the
solution is unique at \(\Sigma_0\), then Proposition~\ref{prop:uniqueness}
applies. This is the standard needed for treating a uniqueness estimator as a
Guttman communality rule. A numerical algorithm may approximate this target, but
the target itself must be exact on the model manifold.

\section{What not to certify}

Several common initial communality guesses are useful starts but are not
consistent estimators of \(c_i\) for a finite congeneric factor block.

\paragraph{Squared multiple correlation.}
For a covariance matrix \(\Sigma\), the SMC communality is
\[
  h_i^{\mathrm{SMC}}
  =
  \sigma_{ii}-\{(\Sigma^{-1})_{ii}\}^{-1}.
\]
In a one-factor block, write
\[
  \Sigma=\Theta+\phi\lambda\lambda\T,\qquad
  t_i=\frac{\phi\lambda_i^2}{\theta_i},\qquad
  \delta=\phi\sum_m\frac{\lambda_m^2}{\theta_m}.
\]
The Sherman-Morrison formula gives
\[
  \{(\Sigma^{-1})_{ii}\}^{-1}
  =
  \theta_i\frac{1+\delta}{1+\delta-t_i},
\]
and hence
\[
  h_i^{\mathrm{SMC}}
  =
  c_i\frac{\delta-t_i}{1+\delta-t_i}.
\]
This is strictly below \(c_i\) at finite block size when \(c_i>0\) and the
other indicators carry finite information. It approaches \(c_i\) only in a
large-block or near-perfect-measurement limit. Therefore SMC is not an exact
population communality rule for the Guttman estimator.

\paragraph{Bounds and repairs.}
Clamps to \([0,\sigma_{ii}]\), wide lavaan bounds, positive-definite repairs, and
ridge corrections can make a finite-sample procedure safer. They are not
population communality estimators unless the repair is inactive at the target.
If a repair is active at \(\Sigma_0\), the estimator has changed target.

\section{Implementation implications}

For a pure Guttman estimator, \texttt{magmaan} should distinguish
communality rules by their population contract:
\begin{itemize}
\item certified rules: single triads, average-of-ratios Spearman, ratio-of-sums
  Spearman, identity linear moments, weighted linear-moment variants with
  nonzero denominators, one-step triad GMM with a declared correlation
  covariance model, and exact uniqueness-first rules;
\item start-value rules only: SMC, arbitrary lower bounds, and repairs that can
  be active at the population target;
\item regularized variants: any rule that replaces a failed denominator or
  singular uniqueness problem by a ridge or projection should be labeled as a
  regularized Guttman variant.
\end{itemize}

The current frontier Guttman estimator uses the average-of-ratios
Spearman-Harman rule through the shared \texttt{theta\_spearman()} helper, and
then subtracts the resulting residual variance from the diagonal. On an exact
simple CFA block this is population-consistent for
\(\diag(\Lambda\Phi\Lambda\T)\). If we add alternatives, the test for admission
should be the propositions above, not historical popularity as a starting value.
For a recommendable closed-form default, the identity linear-moment rule
\eqref{eq:identity-linear-mom} is the least ratio-sensitive member of the
certified family. For an efficiency-oriented variant, use the one-step GMM
rule \eqref{eq:stacked-gmm} with the weight construction documented in the
estimator output.

\begin{thebibliography}{9}

\bibitem{cudeck1991}
Cudeck, R. (1991).
Noniterative factor analysis estimators, with algorithms for subset and
instrumental variable selection.
\emph{Journal of Educational Statistics}, 16, 35--52.
\href{https://doi.org/10.3102/10769986016001035}{doi:10.3102/10769986016001035}.

\bibitem{dhaene2024}
Dhaene, S., and Rosseel, Y. (2024).
An evaluation of non-iterative estimators in confirmatory factor analysis.
\emph{Structural Equation Modeling}, 31(1), 1--13.
\href{https://doi.org/10.1080/10705511.2023.2187285}{doi:10.1080/10705511.2023.2187285}.

\bibitem{harman1976}
Harman, H. H. (1976).
\emph{Modern factor analysis} (3rd ed.).
University of Chicago Press.

\bibitem{iharakano1986}
Ihara, M., and Kano, Y. (1986).
A new estimator of the uniqueness in factor analysis.
\emph{Psychometrika}, 51, 563--566.
\href{https://doi.org/10.1007/BF02295595}{doi:10.1007/BF02295595}.

\bibitem{spearman1904}
Spearman, C. (1904).
``General intelligence,'' objectively determined and measured.
\emph{The American Journal of Psychology}, 15, 201--293.
\href{https://doi.org/10.2307/1412107}{doi:10.2307/1412107}.

\end{thebibliography}

\end{document}
